\section{Instalación y configuración}

\subsection{Preparar el repositorio}
En el equipo anfitrión, clone o actualice el repositorio autorizado de HomePilot:
\begin{lstlisting}[language=bash]
git clone https://github.com/NezuSas/homepilot.git
cd homepilot
git pull --ff-only
\end{lstlisting}

\subsection{Ejecutar el asistente técnico}
El asistente evita que el técnico tenga que editar manualmente la configuración inicial. Ejecute:
\begin{lstlisting}[language=bash]
bash scripts/install-edge-office.sh --wizard
\end{lstlisting}

El asistente solicita, en orden:
\begin{enumerate}[itemsep=0.15cm]
    \item Perfil de instalación.
    \item Acción deseada: desplegar, preparar sin iniciar o diagnóstico.
    \item Autorización de limpieza segura de caché Docker, si corresponde.
    \item Autorización explícita para HACS e integraciones de comunidad, cuando hay Home Assistant.
    \item Confirmación final de la configuración elegida.
\end{enumerate}

\subsection{Desplegar los servicios}
Cuando el resumen sea correcto, confirme el despliegue desde el asistente. Para una actualización posterior controlada se usa:
\begin{lstlisting}[language=bash]
git pull --ff-only
bash scripts/homepilot-maintenance.sh --deploy --yes
\end{lstlisting}

\subsection{Primera cuenta administrativa}
Abra la interfaz local en el navegador:
\begin{lstlisting}[language=bash]
http://<IP-DEL-EQUIPO>:8080
\end{lstlisting}

Con una base de datos nueva y el bootstrap de desarrollo desactivado, HomePilot solicita crear la primera cuenta. Esta cuenta se convierte en administradora de la instalación. El técnico debe entregar la contraseña al responsable designado y pedirle cambiarla en la primera sesión.

\subsection{Conectar Home Assistant cuando aplica}
Para \texttt{bridge\_ha}, abra el Asistente de instalación de HomePilot e ingrese:
\begin{itemize}[itemsep=0.15cm]
    \item La URL accesible desde el contenedor de HomePilot. En una red Docker compartida suele ser \texttt{http://homeassistant:8123}; en una instalación externa use la IP LAN o el DNS local del equipo.
    \item Un token de acceso de larga duración creado desde el perfil administrador de Home Assistant.
\end{itemize}

Pruebe la conexión antes de guardarla. HomePilot cifra y usa el token exclusivamente para el bridge local.

\subsection{HACS y SonoffLAN}
HACS no se instala silenciosamente. Si el proyecto requiere SonoffLAN u otra integración de comunidad, el técnico debe autorizar esta acción desde el asistente y luego completar la configuración dentro de Home Assistant.

\begin{tcolorbox}[title=Importante]
La instalación de HACS y SonoffLAN puede requerir reiniciar Home Assistant y autenticarse con la cuenta del fabricante. Confirme con el cliente antes de modificar una instalación de Home Assistant existente.
\end{tcolorbox}

\subsection{Descubrir e importar dispositivos}
En HomePilot, abra \textit{Sistema > Descubrimiento}. Importe únicamente los dispositivos aprobados para el proyecto. Después asigne cada uno a un hogar y una estancia. HomePilot conserva su propio inventario, aunque la fuente inicial sea Home Assistant.
